\section{Synthesis Slices}
\label{sec:synthesis}



\providecommand{\synslice}[2]{\mathit{SynSlice}\;#1 \mathbin{\blacktriangleleft} #2}
\providecommand{\minsynslice}[2]{\mathit{MinSynSlice}\;#1 \mathbin{\blacktriangleleft} #2}
\providecommand{\exactsynslice}[2]{\mathit{ExactSynSlice}\;#1 \mathbin{\blacktriangleleft} #2}
\providecommand{\isMin}[1]{\mathsf{IsMinimal}\,(#1)}
\providecommand{\joinsyn}{\sqcup}
\providecommand{\meetsyn}{\sqcap}
\providecommand{\pairsyn}{\times}

\noindent The type information of a subexpression in a bidirectional type system is derived either from \textit{within} the term by synthesis, or from the surrounding context. This section defines \textit{synthesis slices} which concern how to explain an expression's \textit{synthesised} type, that is, the type information \textit{internal} to the expression.

\subsection{Synthesis Slices}
\label{sec:syn-slice-def}

Consider a `program' $P = (\Gamma, e)$, that is, typing assumptions and an expression. When the program synthesises a type (in some type variable context $\Delta$)\begin{fn}Type variable contexts are not sliced as the only information such slices would encode is whether a type variable was used, which can be extracted purely syntactically from the resulting slice. Typing assumptions are different because we can \textit{partially} require an assumption.\end{fn}, all of its type information is contained within the synthesis derivation $D \;:\; \synthesis{e}{\tau}$.

Given the derivation $D$, we may \textit{query} the type $\tau$ by taking a slice $\upsilon$ in $\lfloor \tau \rfloor$. A synthesis slice is a `program slice' $\rho = (\gamma, \varsigma)$ in $\lfloor \Gamma , e\rfloor$ which synthesises a type at least as precise as the query $\upsilon$; this corresponds to a highlighted version of the original program. 

Intuitively, $\upsilon$ specifies the part of $\tau$ we wish to explain; for example, querying $\upsilon = \mathit{Int} \Rightarrow \gap$ for $\tau = \mathit{Int} \Rightarrow \mathit{Bool}$ asks ``which parts of the program account for the domain being $\mathit{Int}$?''. The slice might not include \textit{all} the program regions relevant to the query type, but will necessarily provide a consistent subset which suffices to totally \textit{explain} (independently synthesise) the query.

\begin{definition}[Synthesis slice]
\label{def:syn-slice}
A \textit{synthesis slice} of $D : \synthesis{e}{\tau}$ on $\upsilon \in \lfloor \tau \rfloor$, written $\synslice{D}{\upsilon}$, is a pair $(\gamma, \varsigma) \in \lfloor \Gamma, e\rfloor$ such that $\Delta;\gamma \vdash \varsigma\;{\color{BrickRed}\Uparrow}\;\phi$ for some $\phi \sqsupseteq \upsilon$ (with $\phi \in \lfloor \tau \rfloor$ by graduality, \cref{thm:graduality-syn}).
\end{definition}
\begin{figure}[H]
\centering
\begin{tikzcd}[column sep=large, row sep=large]
& \rho \arrow[r, "\sqsubseteq" {sloped}] \arrow[d, "{\color{BrickRed}\Uparrow}"'] & (\Gamma, e) \arrow[d, "{\color{BrickRed}\Uparrow}"] \\
\upsilon \arrow[r, "\sqsubseteq" {sloped}] & \phi \arrow[r, induced, "\sqsubseteq" {sloped}, "\text{(by graduality)}"', dashed] & \tau
\end{tikzcd}
\caption{Diagram of a synthesis slice; $\upsilon \sqsubseteq \phi$ is the validity witness.}
\label{fig:fundamental-triangle}
\end{figure}

Many such synthesis slices exist. In particular, the original program $(\Gamma, e)$ is necessarily a synthesis slice; and for $\synslice{D}{\gap}$, every slice of $(\Gamma, e)$ is a valid synthesis slice. We refer to the $\upsilon \sqsubseteq \phi$ condition as \textit{validity}. It is intentional that this is \textit{not} an exact equality, as demonstrated by the counterexample below:

\begin{counterexamplebox}
%
\begin{counterexample}[Non-Existence of all Exact Synthesis Slices]\label{lem:exact-not-exist}
Consider $x \mathbin{:} \mathbf{*} \Rightarrow \mathbf{*} \vdash (x, x) {\color{BrickRed}\Uparrow} (* \Rightarrow *) \times (* \Rightarrow *)$, and query $\upsilon = (\mathbf{*} \Rightarrow \gap) \times (\gap \Rightarrow \mathbf{*})$. 

Any strictly smaller slice of either $x : * \Rightarrow *$ or $(x, x)$ synthesises a type strictly less precise than $\upsilon$. Yet the original program synthesises $(* \Rightarrow *) \times (* \Rightarrow *)$, strictly more precise than $\upsilon$.
\end{counterexample}
\end{counterexamplebox}

\subsection{Minimality}
\label{sec:minimality}


Synthesis slices give us \textit{some} explanation of the query, but many are clearly useless, such as the original program itself, which corresponds to fully highlighting the program. We want to find minimal slices: those where further slicing any part of the program makes the slice invalid (synthesising a type insufficient to cover the query).
\begin{definition}[Minimality]
\label{def:isminimal}
For an element $a$ of a poset, $a$ is minimal, $\isMin{a}$, when every $a'$ with $a' \sqsubseteq a$ is in fact equal to $a$.
\end{definition}
\begin{definition}[Minimal Synthesis Slices]
A $\minsynslice{D}{\upsilon}$ is a $\synslice{D}{\upsilon}$ minimal in $\synslice{D}{\upsilon}$ under program-slice precision on $\lfloor \Gamma, e\rfloor$.
\end{definition}

\subsection{Existence of Minimal Slices}
\label{sec:min-exists}


Crucially, any synthesis slice has a minimal slice below it (or is minimal):
\begin{majortheorembox}
%
\begin{theorem}[Existence of minimal slices]\label{thm:min-exists}
For every $s : \synslice{D}{\upsilon}$, there exists $m : \minsynslice{D}{\upsilon}$ with $m \sqsubseteq s$.
\end{theorem}
\end{majortheorembox}

\begin{proof}
The slice lattice $\lfloor \Gamma, e \rfloor$ is finite. Hence, $\{\,s' : \synslice{D}{\upsilon} \mid s' \sqsubset s\,\}$ is finite and computable by enumeration. If this set is empty, then $s$ is minimal; otherwise, recurse on any of its elements. Each step strictly decreases precision, so finiteness ensures termination.
\end{proof}

Consequentially, any derivation $D$ has a minimal synthesis slice, as the original program is a trivial inhabitant of $\synslice{D}{\upsilon}$ for all $\upsilon$.
\paragraph{A Brute-Force Algorithm via Graduality} The proof above is not a practical algorithm. In reality, for any term, you can efficiently compute the list of maximal slices that are strictly less precise than some given term by omitting exactly \textit{one} leaf from the AST. 

Pick the first maximal strict slice that satisfies the query and recurse; if none of the slices satisfy the query, then you have found a minimal slice by using graduality (\cref{thm:graduality-syn}). This is $\mathcal{O}(n \times T)$ where $n$ is the program size and $T$ the type checking complexity; type checking is linear in the program size for this calculus, so finding a minimal slice is $\mathcal{O}(n^2)$.

See Figure~\ref{fig:syn-slice-existence} for an example lattice to descend through, stopping just before descending below the given query; all three minima are marked.

\tikzset{
  slprog/.style={draw=latticecolor!55, fill=white, rounded corners=2pt,
                 inner sep=2.2pt, font=\scriptsize},
  slmin/.style={draw=OliveGreen!70!black, fill=white, line width=0.9pt,
                ellipse, minimum height=0.95cm, inner xsep=2pt, inner ysep=2pt, font=\scriptsize},
  slexact/.style={draw=OliveGreen!75!black, fill=white, line width=1.2pt,
                  rectangle, minimum width=0.85cm, minimum height=0.85cm, inner sep=2pt, font=\scriptsize},
  slhi/.style={draw=Green!60!black, fill=white, line width=1.3pt, rounded corners=2pt,
               inner sep=2.2pt, font=\scriptsize},
  slprogarr/.style={->, latticecolor!75, line width=0.95pt},
  slsyn/.style={->, BrickRed, line width=0.6pt, dashed,
                dash pattern=on 3pt off 2pt, shorten >=2.5pt, shorten <=2.5pt},
  slpath/.style={->, Green!65!black, line width=1.5pt, shorten >=2pt, shorten <=2pt},
  slfront/.style={OliveGreen!55!black, line width=1.5pt, line cap=round},
  sldot/.style={OliveGreen!55!black, line width=0.5pt, dotted},
}
\newcommand{\sliceNodes}{%
  \node[slprog] (T)   at ( 0.0, 9.0) {$(P,\;\tau)$};
  \node[slprog] (A)   at (-2.7, 7.4) {$(\rho_a,\;\phi_a)$};
  \node[slprog] (B)   at ( 0.1, 7.3) {$(\rho_b,\;\phi_b)$};
  \node[slprog] (C)   at ( 2.7, 7.4) {$(\rho_c,\;\phi_c)$};
  \node[slprog] (E)   at (-1.9, 5.9) {$(\rho_e,\;\phi_e)$};
  \node[slprog] (F)   at ( 1.4, 6.0) {$(\rho_f,\;\phi_f)$};
  \node (G)   at ( 3.0, 5.0) {};
  \node[slprog] (H)   at (-3.0, 4.4) {$(\rho_h,\;\phi_h)$};
  \node[slprog] (I)   at (-0.5, 4.2) {$(\rho_i,\;\upsilon)$};
  \node[slprog] (J)   at ( 2.5, 3.7) {$(\rho_j,\;\phi_j)$};
  \node[slprog] (K)   at (-1.7, 2.5) {$(\rho_k,\;\phi_k)$};
  \node[slprog] (Bot) at ( 0.7, 1.1) {$(\gap,\;\gap)$};
}
\newcommand{\sliceEdges}{%
  \foreach \p/\c in {T/A,T/B,T/C,A/E,A/H,B/E,B/F,C/F,E/I,F/I,F/J,H/K,I/K,I/Bot,J/Bot,K/Bot}
    \draw[slprogarr] (\p) -- (\c);
}
\newcommand{\sliceFrontier}{%
  \draw[slfront] (-3.9,3.5) .. controls (-2.4,3.5) and (-1.1,4.7) .. (3.9,4.3);
}

\begin{figure}[t]
\centering
\begin{subfigure}[t]{0.49\linewidth}
\centering
\begin{tikzpicture}[scale=0.72, transform shape, >={Stealth[length=4pt,width=3pt]}]
\begin{scope}[on background layer]
  \fill[OliveGreen!14] (-3.9,9.8) -- (-3.9,3.5)
    .. controls (-2.4,3.5) and (-1.1,4.7) .. (3.9,4.3) -- (3.9,9.8) -- cycle;
  \sliceFrontier
  \draw[sldot] (-3.9,9.8) -- (3.9,9.8)  (-3.9,9.8) -- (-3.9,3.5)  (3.9,9.8) -- (3.9,4.3);
\end{scope}
\sliceNodes
\sliceEdges
\draw[slprogarr] (G) -- (J);
\draw[slprogarr] (C) -- (G);
\foreach \p/\c in {T/A,T/B,T/C,A/E,A/H,B/E,B/F,C/F,C/G,E/I,F/I,F/J,G/J,H/K,I/K,I/Bot,J/Bot,K/Bot}
  \draw[slsyn] (\p) to[bend right=7] (\c);
\draw[slsyn] (F) -- (G);
\draw[slsyn] (E) -- (H);
\node[slmin] at (G) {$(\rho_g,\;\phi_g)$};
\node[slmin] at (H) {$(\rho_h,\;\phi_h)$};
\node[slexact] at (I) {$(\rho_i,\;\upsilon)$};
\node[anchor=north west, font=\scriptsize\bfseries, OliveGreen!35!black] at (-3.85,9.6) {$\upsilon \sqsubseteq \phi_i$};
\node[anchor=south west, font=\scriptsize\bfseries, gray!50!black] at (-3.85,0.7) {$\upsilon \not\sqsubseteq \phi_i$};
\end{tikzpicture}
\caption{Valid synthesis slices above query $\upsilon$ in the shaded region, with 3 minima consisting of 2 inexact minima (ellipses) and one minimum (square) synthesising exactly $\upsilon$.}
\label{fig:syn-slice-existence}
\end{subfigure}\hfill
\begin{subfigure}[t]{0.49\linewidth}
\centering
\begin{tikzpicture}[scale=0.72, transform shape, >={Stealth[length=4pt,width=3pt]}]
\begin{scope}[on background layer]
  \fill[OliveGreen!8]  (-3.9,9.8) -- (-3.9,2.2) .. controls (-1.2,1.8) and (1.2,1.8) .. (3.9,2.2) -- (3.9,9.8) -- cycle;
  \fill[OliveGreen!18] (-3.9,9.8) -- (-3.9,3.5) .. controls (-2.4,3.5) and (-1.1,4.7) .. (3.9,4.3) -- (3.9,9.8) -- cycle;
  \fill[OliveGreen!32] (-3.9,9.8) -- (-3.9,6.6) .. controls (-1.5,6.9) and (1.5,7.0) .. (3.9,6.7) -- (3.9,9.8) -- cycle;
  \draw[slfront] (-3.9,6.6) .. controls (-1.5,6.9) and (1.5,7.0) .. (3.9,6.7);
  \sliceFrontier
  \draw[slfront] (-3.9,2.2) .. controls (-1.2,1.8) and (1.2,1.8) .. (3.9,2.2);
  \draw[sldot] (-3.9,9.8) -- (3.9,9.8)  (-3.9,9.8) -- (-3.9,2.2)  (3.9,9.8) -- (3.9,2.2);
\end{scope}
\sliceNodes
\sliceEdges
\foreach \p/\c in {T/A,T/B,T/C,A/E,A/H,B/E,B/F,C/F,E/I,F/I,F/J,H/K,I/K,I/Bot,J/Bot,K/Bot}
  \draw[slsyn] (\p) to[bend right=7] (\c);
\draw[slsyn] (E) -- (H);
\draw[slpath] (T) -- (B);
\draw[slpath] (B) -- (E);
\draw[slpath] (E) -- (I);
\draw[slpath] (I) -- (K);
\draw[slpath] (K) -- (Bot);
\node[slhi] at (B) {$(\rho_b,\;\phi_b)$};
\node[slhi] at (I) {$(\rho_i,\;\upsilon_1)$};
\node[slhi] at (K) {$(\rho_k,\;\phi_k)$};
\node[anchor=west, font=\scriptsize\bfseries, Green!50!black] at ([xshift=3pt]B.east) {$m_2$};
\node[anchor=east, font=\scriptsize\bfseries, Green!50!black] at ([xshift=-3pt]I.west) {$m_1$};
\node[anchor=east, font=\scriptsize\bfseries, Green!50!black] at ([xshift=-3pt]K.west) {$m_0$};
\node[anchor=west, font=\scriptsize, OliveGreen!40!black] at (3.95,6.7) {$\upsilon_2$};
\node[anchor=west, font=\scriptsize, OliveGreen!40!black] at (3.95,4.3) {$\upsilon_1$};
\node[anchor=west, font=\scriptsize, OliveGreen!40!black] at (3.95,2.2) {$\upsilon_0$};
\end{tikzpicture}
\caption{Monotonically descending minima $m_2, m_1, m_0$ for query refinements $\upsilon_2 \to \upsilon_1 \to \upsilon_0$, calculated by descending along the \textcolor{Green!60!black}{green} path.}
\label{fig:syn-slice-refinement}
\end{subfigure}
\caption{Two similar Hasse diagrams of program slices $\rho_i \in \lfloor P \rfloor$, ordered by {\color{latticecolor}strict program precision $\sqsubset$}, paired with corresponding synthesised types $\phi_i \in \lfloor \tau \rfloor$. Overlaid with precision relations on types ({\color{BrickRed}dotted lines}).}
\label{fig:syn-slices}
\end{figure}

\subsection{Refined Slices}
\label{sec:sub-slices}
In practice, the user may look at a large query $\upsilon$ getting a large program slice $\rho$. Therefore they may want to refine the query to focus on a subpart, $\upsilon' \sqsubseteq \upsilon$, that they don't quite understand. It turns out to always be possible to pick a consistent program slice $\rho' \sqsubseteq \rho$ which is a refinement of the previous slice by \textit{monotonicity} (which follows from graduality, \cref{thm:graduality-syn}):

\begin{theorem}[Monotonicity of minimal slices]\label{thm:mono}
If $\upsilon_1 \sqsubseteq \upsilon_2$ in $\lfloor\tau\rfloor$ and $m_2 : \minsynslice{D}{\upsilon_2}$, then there exists $m_1 : \minsynslice{D}{\upsilon_1}$ with $m_1 \sqsubseteq m_2$.
\end{theorem}

This is desirable because the user will not be shown \textit{additional} information after a query refinement which would require further reasoning by the user to understand. \Cref{fig:syn-slice-refinement} illustrates a monotonic refinement path along 3 refined queries on the example lattice. If we did not follow a monotonic path during refinement we would get `branch hopping':

\paragraph{Branch hopping.} This is where highlighted regions switch between branches when refining the query, requiring the user to switch their reasoning between the contents in each branch. For example, refining $\upsilon = \mathbf{*}\Rightarrow\mathbf{*}$ to $\upsilon' = \gap\Rightarrow\mathbf{*}$:
\begin{center}
\fullslice{x : \sli{\mathbf{*} \Rightarrow \mathbf{*}}}{\sli{\mathbf{if}}\;\mathbf{true}\;\sli{\mathbf{then}\;x\;\mathbf{else}}\;y}
\quad$\xrightarrow{\upsilon' \sqsubseteq \upsilon}$\quad
\fullslice{y : \mathbf{*}\,\sli{\Rightarrow \mathbf{*}}}{\sli{\mathbf{if}}\;\mathbf{true}\;\sli{\mathbf{then}}\;x\;\sli{\mathbf{else}\;y}}
\end{center}

On the other hand, you do not necessarily have monotonicity in the upwards direction. Consider the $\rho_j$ node in \Cref{fig:syn-slice-refinement}, which is minimal under query $\upsilon_0$, but no minimum for query $\upsilon_1$ is strictly more precise than $\rho_j$.

\subsection{Decompositions}
\label{sec:decompositions}
\mastercases

Instead of calculating slices by brute force (\cref{sec:min-exists}), it would be ideal if we could compositionally calculate minimal slices of a term from minimal subslices of its subterms. It is indeed possible to compose synthesis slices into larger synthesis slices, and we prove that all minima can be expressed by some composition; this subsection demonstrates two of the cases. \Cref{sec:algorithms} considers how to recursively choose the compositions.

\subsubsection{Products}
\label{sec:case-pair}

Given slices $(\gamma_1, \varsigma_1) : \synslice{D_1}{\upsilon_1}$ and $(\gamma_2, \varsigma_2) : \synslice{D_2}{\upsilon_2}$ for the two components of a pair $D = \synr{Pair}\,D_1\,D_2$ (with synthesised types $\phi_1, \phi_2$), we form their product slice $(\gamma_1, \varsigma_1) \pairsyn (\gamma_2, \varsigma_2) : \synslice{D}{\upsilon_1 \times \upsilon_2}$ by joining the two assumption slices and pairing the term slices:
\[(\gamma_1, \varsigma_1) \pairsyn (\gamma_2, \varsigma_2) \;\triangleq\; (\gamma_1 \sqcup \gamma_2,\;(\varsigma_1, \varsigma_2))\]
The pair synthesises (possibly more precise)\begin{fn}Due to $\gamma_1 \sqcup \gamma_2$ being larger than $\gamma_1$ and $\gamma_2$\end{fn} types $\phi_1' \times \phi_2' \sqsupseteq \phi_1 \times \phi_2 \sqsupseteq \upsilon_1 \times \upsilon_2$. For this combinator, we get the following decomposition theorem:

\begin{thmcase}[Product Decomposition]\label{thm:min-prod-decomp}
Let $m : \minsynslice{\synr{Pair}\,D_1\,D_2}{\upsilon_1 \times \upsilon_2}$. Then there exist minimal $m_1 : \minsynslice{D_1}{\upsilon_1}$ and $m_2 : \minsynslice{D_2}{\upsilon_2}$ with $m = m_1 \pairsyn m_2$.
\end{thmcase}

However, we shall later see that the converse fails: the product of two minimal slices need not be minimal, since the joined assumptions can leave one component's term slice redundant. So one cannot simply recurse through the derivation, pairing sub-slices on $\upsilon_1$ and $\upsilon_2$ to obtain a minimal slice on $\upsilon_1 \times \upsilon_2$. \Cref{sec:product-revisited} explores and resolves this problem.

\subsubsection{Let Bindings}
\label{sec:let-decomp}

A more subtle case is that of let-bindings, which have different assumptions for their two premises:
\[\inference[\synr{Let}]{
  \synthesis{e_1}{\tau_1}
  &
  \synthesis[\Delta;\,\Gamma, x:\tau_1]{e_2}{\tau_2}
}{\synthesis{\mathbf{let}\;x = e_1\;\mathbf{in}\;e_2}{\tau_2}}\]
A slice of a let-expression can be constructed from a slice $(\gamma_1, \varsigma_1)$ of the definition (synthesising $\phi_1$) and a slice $(\gamma_2, \varsigma_2)$ of the body, constrained with an assumption consistency condition: the body's used assumptions for the bound variable $x$ must be \textit{at most as precise as} the type synthesised by the slice of the definition, $\gamma_2(x) \mathbin{\sqsubseteq} \phi_1$.

Then, upon joining the assumptions for the two slices, we must remove the assumption on $x$ from the body (notated $\gamma_{\setminus x}$), as this is now provided by the definition:
\[(\gamma_1, \varsigma_1) \mathbin{\mathsf{def}} (\gamma_2, \varsigma_2) \triangleq (\gamma_1 \sqcup (\gamma_2)_{\setminus x},\; \mathbf{let}\;x = \varsigma_1\;\mathbf{in}\;\varsigma_2)\]

\begin{thmcase}[Binding Decomposition]\label{thm:min-def-decomp}
Let $m : \minsynslice{\synr{Let}\,D_1\,D_2}{\upsilon}$. Then there exists $\phi_1 \in \lfloor\tau_1\rfloor$ and minimal slices $m_1 : \minsynslice{D_1}{\phi_1}$ (synthesising $\phi_1'$) and $m_2 = (\gamma_2, \varsigma_2) : \minsynslice{D_2}{\upsilon}$ satisfying $\gamma_2(x) \sqsubseteq \phi_1'$, with $m = m_1 \mathbin{\mathsf{def}} m_2$.
\end{thmcase}

\subsection{Contribution Slices}
\label{sec:join-slices}


Finally, a minimal slice only provides an explanation; it does not show \textit{all} information \textit{contributing} in any way to a query type. The latter is a useful notion when performing refactors which change the type, where quickly identifying the regions of the program which rely on the old type is crucial.

Such a region can be uniquely expressed as the \textit{join} of all minimal slices for the query. The contribution slice is efficient in the sense that it is a \textit{least} upper bound of actual minimal slices. For this to work, synthesis slices must be closed under joins: the joined slice must also be \textit{valid}, which is true:

\begin{theorem}[Join Closure of Synthesis Slices]\label{thm:join-syn}
For $(\gamma_1, \varsigma_1) : \synslice{D}{\upsilon_1}$ and $(\gamma_2, \varsigma_2) : \synslice{D}{\upsilon_2}$, then $(\gamma_1, \varsigma_1) \joinsyn (\gamma_2, \varsigma_2) : \synslice{D}{\upsilon_1 \sqcup \upsilon_2}$, with the joined slice taken pointwise: $(\gamma_1 \sqcup \gamma_2,\; \varsigma_1 \sqcup \varsigma_2)$.
\end{theorem}

\begin{figure}[H]
\centering
\begin{tikzcd}[column sep=large, row sep=large]
\rho_1 \arrow[r, "\sqsubseteq"] \arrow[d, "{\color{BrickRed}\Uparrow}"'] & {\color{induced} \rho^{\sqcup}} \arrow[d, "{\color{BrickRed}\Uparrow}"] & \rho_2 \arrow[l, "\sqsupseteq"'] \arrow[d, "{\color{BrickRed}\Uparrow}"] \\
\phi_1 & {\color{induced} \phi'} & \phi_2 \\
\upsilon_1 \arrow[u, "\sqsubseteq" {sloped}] \arrow[r, "\sqsubseteq"] & {\color{induced} \upsilon_1 \sqcup \upsilon_2} \arrow[u, induced, "\sqsubseteq"' {sloped}, dashed] & \upsilon_2 \arrow[u, "\sqsubseteq"' {sloped}] \arrow[l, "\sqsupseteq"']
\end{tikzcd}
\caption{Theorem~\ref{thm:join-syn}.}
\label{fig:join-validity}
\end{figure}

%%% Local Variables:
%%% mode: LaTeX
%%% TeX-master: "PolymorphicTypeSlicing"
%%% End:
